% !TeX root = ../nc-quadric-surface-avatar.tex
% 2026-11-23 / v1.0
\section{相対heartの比較とcubic代数族の回復}
\label{sec:qa-relative-avatar-heart-comparison}

\subsection{相互の生成対象による相対 heart の認識}
\label{subsec:qa-relative-mutual-heart-recognition}

次の議論は fiberwise な $t$-exactness の持上げではない．
両側で既に構成された相対 $H_2$ heart の中に，互いの構造対象が
実際に存在することを用いる．したがって非被約基底でも同じ証明が成立する．

\begin{lemma}[生成族に対する非正性の検出]
\label{lem:qa-generators-detect-right-texactness}
$\mathcal A,\mathcal B$ を局所 noetherian Grothendieck 圏とし，
$\mathcal A$ の noetherian 部分を $\mathcal A^{\rm noeth}$ とする．
$\{L_\alpha\}$ を $\mathcal A^{\rm noeth}$ 内の族であって，
全ての noetherian 対象がその有限直和の商であると仮定する．
標準の安定導来 enhancement の間の colimit を保つ完全関手
\[
 \mathsf F:D(\mathcal A)\longrightarrow D(\mathcal B)
\]
が全ての $L_\alpha$ を $D^{\leq0}(\mathcal B)$ に送るなら，
$\mathsf F$ は右 $t$-exact，すなわち
$\mathsf F(D^{\leq0}(\mathcal A))\subset D^{\leq0}(\mathcal B)$ である．
\end{lemma}
\begin{proof}
任意の noetherian 対象 $M$ に対し，生成族からの全射を取り，
その kernel にも同じ操作を繰り返す．kernel も noetherian なので，
\[
 \cdots\longrightarrow Q^{-2}\longrightarrow Q^{-1}
       \longrightarrow Q^0\longrightarrow M\longrightarrow0
\]
という完全列を得る．各 $Q^{-j}$ は $L_\alpha$ の有限直和である．
これは射影解消である必要はなく，通常の複体 $Q^\bullet$ が
$M$ とquasi-isomorphicであることだけを使う．

$Q^\bullet$ は，区間 $[-n,0]$ に制限した brutal truncation
$\sigma^{\geq-n}Q^\bullet$ の順方向 homotopy colimit である．
実際，これらは複体の包含写像による系であり，filtered colimit の
完全性からその homotopy colimit は項別 colimit とquasi-isomorphicである．
各有限複体の像は $D^{\leq0}(\mathcal B)$ に属する．
同部分圏は coproduct と extension に閉じているので，
sequential homotopy colimit にも閉じている．したがって
$\mathsf F(M)\in D^{\leq0}(\mathcal B)$ である．

任意の $\mathcal A$ の対象は noetherian 部分対象の filtered colimit である．
導来圏の標準 enhancement において，この colimit は homotopy colimit
とも一致し，$D^{\leq0}$ はこれらのfiltered homotopy colimitに閉じている．
これはGrothendieck圏のAB5による．$\mathsf F$ はこれを保つので，任意の heart 対象も
非正部分に送られる．
最後に，$D^{\leq0}(\mathcal A)$ の任意の対象を次数 $0$ 以下の複体で表し，
再び有限 brutal truncation の homotopy colimit として書けば結論を得る．
本稿で用いる $\mathsf F$ は perfect Morita 関手の Ind 延長であり，
この enhancement と全 colimit の保存を備えている．
\end{proof}

\begin{proposition}[相対 $H_2$ heart の相互標識判定]
\label{prop:qa-h2-mutual-markings-heart-test}
同じ基底 $S$ 上の二つの相対 $H_2$ 族を
$(\mathcal C,\mathcal O)$，$(\mathcal D,\mathcal P)$ とする．
共通の fppf 局所 chart 上の compact 圏の間に，
任意基底変更と整合する $S$-線形 Morita 同値
\[
 \mathsf F^c:\mathcal D^{\rm perf}\xrightarrow{\sim}
                  \mathcal C^{\rm perf}
\]
があるとする．さらに，実際の相対 heart に関する条件
\[
 \mathsf F^c(\mathcal P)\in\mathcal C,
 \qquad
 (\mathsf F^c)^{-1}(\mathcal O)\in\mathcal D
 \tag{MH}\label{eq:qa-mutual-heart-condition}
\]
がその chart の全ての affine 基底変更上で成立すると仮定する．
このとき $\mathsf F^c$ は唯一の colimit を保つ導来 Morita 同値へ延長され，
その延長は $t$-exact である．従って基底変更と整合する完全同値
\[
 \mathsf F:\mathcal D\xrightarrow{\sim}\mathcal C
\]
を得る．初めの compact 比較に coherent な降下データがある場合，
この完全同値にも同じ降下データがある．
\end{proposition}
\begin{proof}
相対 $H_2$ の定義により，各局所 affine chart 上で
$D(\operatorname{Ind}\mathcal C)$ と
$D(\operatorname{Ind}\mathcal D)$ は compactly generated である．
従って perfect Morita 同値はこれらの間の colimit を保つ同値に
一意に延長される（\cite[\S4.3,\S8.2]{Keller1994}の標準dg Morita）．

canonical な相対反 Serre 関手をそれぞれ $t_{\mathcal C},t_{\mathcal D}$
と書く．両者は heart 上で完全であり，Serre 双対の一意性により
\[
 \mathsf F^c t_{\mathcal D}\simeq
 t_{\mathcal C}\mathsf F^c
\]
が成り立つ．論文Bの
\cite[\texttt{prop:relative-h2-antiserre-reconstruction}]{QuadricModuli2026}
により，$\{t_{\mathcal C}^{-n}\mathcal O\}_{n\geq0}$ と
$\{t_{\mathcal D}^{-n}\mathcal P\}_{n\geq0}$ はそれぞれ AZ 生成族である．
条件~\eqref{eq:qa-mutual-heart-condition} と $t$ の完全性から，
これらの生成族は $\mathsf F$ とその逆によって相手の heart に送られる．
前補題を両向きに適用すると，$\mathsf F$ と $\mathsf F^{-1}$ は
ともに右 $t$-exact である．従って
$\mathsf F(D^{\leq0})=D^{\leq0}$ であり，直交性から非負部分も保つ．
これで $t$-exactness が従う．

得られる Ind heart の完全同値は noetherian 対象を保つので，
指定された coherent heart の完全同値に制限される．
任意の基底変更に対する導来比較は compact 生成対象上の比較から
一意に延長される．通常の module-object 引戻しはその次数零 cohomology
なので，$t$-exactness によって通常の引戻しとの coherent な比較を得る．
降下データの間の自然同型も compact 生成対象から一意に延長される．
したがって局所判定を貼り合わせられる．ここでは，任意の相対 compact
圏に独立の heart を与えたり，fiberwise な cohomology 一致から heart を
復元したりしていない．
\end{proof}

\begin{corollary}[中間頂点も同じ相対 heart 内で認識される]
\label{cor:qa-relative-middle-same-heart}
定理~\ref{thm:qa-relative-avatar-middle-h2-realizations} の記法で，
その compact Morita 比較は完全同値
\[
 \mathcal C_i\xrightarrow{\sim}\mathcal C_0|_{U_i},
 \qquad \mathcal O_i\longmapsto E_i\quad(i=1,2)
\]
へ延長される．従って，元の同じ相対 heart の対
$(\mathcal C_0|_{U_i},E_i)$ が実際に $\mathfrak H_2(U_i)$ の対象である．
この主張と比較は任意基底変更と整合する．
\end{corollary}
\begin{proof}
前節では全 $E_j$ が $\mathcal C_0$ の基底平坦な heart 対象であることを
既に示した．従って $\mathsf F^c(\mathcal O_i)=E_i$ は第一条件を満たす．
一方，$\mathcal N_i$ の第四対象は $t_\Lambda E_0$ である．
真正な normal foundation の第四対象は $\mathcal C_i$ の基底平坦な
heart 対象だから，canonical $t_i^{-1}$ を適用して
\[
 (\mathsf F^c)^{-1}(E_0)
   =t_i^{-1}(\text{$\mathcal N_i$ の第四対象})\in\mathcal C_i
\]
を得る．$t_i^{-1}$ は完全であり基底変更と可換するので，
この対象も基底上平坦であり，全ての affine 基底変更後に同じ議論が成立する．
これが第二条件である．前命題を適用すればよい．
頂点の相対枠を変えた際の基底線束との tensor は，heart の完全自己同値であり，
この二条件と基底平坦性を保つ．従って実際の cone に含まれる determinant
線束や，剛体化した foundation chart の頂点線束を自明化する必要もない．
\end{proof}

\begin{remark}[この判定の範囲]
前命題で必要なのは互いの構造対象が相手の実際の heart に属するという
二つの条件である．その検証を，幾何 fiber で heart に属することだけで
置き換えることはしない．また相手が未構成の一般 AS 商圏である場合には，
compact generation と相対 AZ 生成族の存在が別途必要である．
したがって本系は，任意の相対 AS/QP stack と $H_2$ stack の同一視ではない．
\end{remark}

\begin{theorem}[同じ相対heartの全標準構造頂点の開条件]
\label{thm:qa-relative-all-vertex-h2-open}
前節のflat曲線型square avatarから構成した相対heart $\mathcal C_0$ と
その四つの基底平坦な対象 $E_i$ を取る．
局所有限型siteの任意の底変更 $S'\to S$ について，
\[
 (\mathcal C_{0,S'},E_{i,S'})\in\HHstack(S')
 \quad\Longleftrightarrow\quad S'\to S\text{ が }U_i\text{ を経由する},
\]
ただし
\[
 U_0=U_3=S,\qquad U_1=D(\det\mu_F),\qquad U_2=D(\det\mu_G)
\]
である．したがってこの二行列式開は，幾何fiberの適合性だけでなく，
\emph{同じ実際の相対heart}の中の強い$\HH$構造対象条件を表す．
必要な根・AS4被覆は局所存在条件のままであり，標識へ加えない．
\end{theorem}
\begin{proof}
$U_0,U_3$ では定理~\ref{thm:qa-relative-square-e0-e3-h2}，
$U_1,U_2$ では前系を用いる．
いずれも任意のsite内の底変更と整合する．
逆に相対$\HH$の条件は全幾何fiberの$\HH$所属を含むので，
定理~\ref{thm:qa-relative-vertex-determinant-open}によって
底変更は該当する開を経由する．
\end{proof}


\subsection{有限窓からの実際のcubic族}

\begin{proposition}[avatarのnormal窓からの相対cubic代数族]
\label{prop:qa-relative-cubic-from-avatar-window}
flat な曲線型 square avatar の族に対し，
$U_2=D(\det\mu_G)$ 上の normal 窓
\[
 (E_0,E_2,E_3,Z_1),\qquad
 Z_1=\operatorname{Cone}
   (E_1\to\mathcal H_{13}^{\vee}\otimes E_3)
\]
を取る．その実際の有限 End 代数を $\Gamma_F$ とする．
すると非線型 geometric quintuple の相対族が基底変更と可換に復元され，
その cubic $\mathbb Z$-代数族の初期四項の End は
$\Gamma_F$ と全ての積を含めて同型になる．
全次数の成分は有限局所自由で，標準 cubic 解消と AS 双対は
任意基底変更後にも成立する．
局所 affine な Noetherian 底では全 projective row と column は
左右 Noetherian であり，任意の Noetherian 基底変更後も同様である．

この cubic 族の canonical tails は，実際の相対 $H_2$ 族である．
その指定零対象からの normal foundation は $\lambda_F$ を与えるので，
$\Xi^{-1}\lambda_F$ によって構成した $E_0$ の相対 $H_2$ 族と
pointed exact 同値である．
$U_1=D(\det\mu_F)$ 上にも $F,G$ を交換した同じ主張が成立する．
\end{proposition}
\begin{proof}
\emph{(1) 有限ベクトル束から実際のquintupleを作る．}
記号を
\[
 \mathcal U=\mathcal H_{02},\qquad
 \mathcal V=\mathcal H_{23},\qquad
 \mathcal W=\mathcal H_{13}^{\vee},\qquad
 C=\operatorname{coker}
   (\mathcal H_{01}\to\mathcal H_{13}^{\vee}\otimes\mathcal H_{03})
\]
とする．三つの束 $\mathcal U,\mathcal V,\mathcal W$ は階数二，
$C$ は階数六である．$U_2$ 上では第一の合成
\[
 \mathcal V\otimes\mathcal U\xrightarrow{\sim}\mathcal H_{03}
\]
が同型であり，最後の合成は全射なので
\[
 0\to R\longrightarrow
 \mathcal W\otimes\mathcal V\otimes\mathcal U
 \longrightarrow C\to0
\]
という階数二の核を持つ局所分裂完全列を得る．
従って包含をテンソルに読み替えると
\[
 w_R\in
 \mathcal U\otimes\mathcal V\otimes\mathcal W\otimes R^\vee
\]
となる．これは nowhere zero で，四つの階数二の束と
可逆なテンソル直線 $\OO\,w_R$ を与える．
全構成はベクトル束の核・双対・tensor と実際の積から来るので，
非平坦なものを含む任意基底変更と可換する．

これは論文B \cite[\texttt{prop:cubic-foundation-open}]{QuadricModuli2026} の構成そのものである．
各幾何 fiber は第一mutationで得た cubic AS 代数の有限窓だから，
$w_R$ は各 fiber で geometric である．
geometric 条件は双対射影束上の proper incidence の補集合という
実際の開条件なので，この族全体が geometric 開へ因数分解する．
枠を選ぶ局所化を除けば，この操作は avatar の全同型にも自然である．

\emph{(2) linear fiber がないこと．}
もしある幾何 fiber が linear cubic であれば，その標準四項を
\[
 (H_0,H_1,H_2,H_3)
 =(\OO,\OO(1,0),\OO(1,1),\OO(2,1))
\]
と取れる．逆mutationの余評価列は通常の pencil 評価列であり，
\[
 M=\operatorname{coker}(\OO\to\OO(1,0)^{\oplus2})
   =\OO(2,0).
\]
従って逆square窓 $(H_1,M,H_2,H_3)$ は，
$\OO(-1,0)$ を共通tensorすると
\[
 (\OO,\OO(1,0),\OO(0,1),\OO(1,1))
\]
という classical square になる．
補題~\ref{lem:qa-as10-all-index-mutual-inverse} の全整数での互逆性と
定理~\ref{thm:qa-as-avatar-bijection} の同型類分類により，
元の square avatar も classical でなければならない．
これは入力が Cartier 曲線型であることに反する．
有限 End から回復される一点表現も，曲線型では一次元の $D$，
classical では二次元の $Y$ なので両型は区別される．
従って全族は非線型 geometric 開へ因数分解する．
幾何 fiber に linear 点がないことからこの\emph{開}への因子化を
結論しているので，非被約基底でも問題はない．

\emph{(3) 全次数とcanonical tails．}
復元したquintupleから，三次関係を $w_R$ の縮約として定め，
四因子を巡回して全整数の cubic $\mathbb Z$-代数を作る．
初期四頂点では二段階合成に関係がなく，三段階合成の関係は
正に上の $R$ である．従ってその有限 End は $\Gamma_F$ と一致する．
この同定は単に Hom 階数の一致ではなく，全ての積を保持する．

有限型の非線型 geometric 普遍族の
論文B \cite[\texttt{lem:universal-flatness}]{QuadricModuli2026} を引き戻せば，
全次数の局所自由性，標準 cubic 解消，AS 双対，および全基底変更との
整合性を得る．非被約な入力基底の行列の fiber 階数一定性だけから
平坦性を推論する必要はない．
同論文 \cite[\texttt{thm:relative-geometric-cubic-noetherian}]{QuadricModuli2026} が
任意 Noetherian 基底変更後の row/column Noether 性を与え，
\cite[\texttt{prop:general-cubic-canonical-tails}]{QuadricModuli2026} が実際の module-object
基底変更と整合する canonical tails，四項 tilting と Serre 比較を与える．
右加群の添字反転は
定理~\ref{thm:qa-geometric-cubic-relative-forgetting} の規約に従う．

同定した初期四項を $Q_0,Q_1,Q_2,Q_3$ とし，
forward の添字規約で書けば
\[
 (Q_0,Q_1,Q_2,Q_3)=(E_0,E_2,E_3,Z_1)
\]
である．有限 End の相対 Morita 同値は相対 Serre functor を保つので，
canonical inverse dualizing complex とその各非負 tensor 冪は
$\Gamma_F$ のものと一致し，この比較は積と基底変更に整合する．
ここで得た実際の cubic 代数族に，四次 overlap の可逆束や
頂点線束を省略した厳密な四周期を追加で課すことはしない．
全整数添字付き代数と canonical Serre 完成の\emph{厳密な}表示同一視には，
これらの可逆束を含む vertex-line transport の明記が別に必要である．

\emph{(4) 既存の相対H2射とnormal窓の認識．}
定理~\ref{thm:qa-geometric-cubic-relative-forgetting} を適用すれば，
上の canonical tails 自体が $Q_0$ を構造対象とする
実際の $\HHstack$ 族となる．
非線型の特異・可約・非被約曲線もこの定理の範囲に入る．

その標準四項が相対的にも真正な normal 窓であることを確認する．
適合二周期を選ぶ fppf 局所では $Q_2=sQ_0$ が実際に成立する．
真の局所 normal ruling を $X,Y$ とすると，$Q_1$ は
各幾何 fiber で $X$ または $Y$ の一方である
（論文A \cite[\texttt{thm:cubic-coverage}]{QuadricIntrinsic2026} の必要性証明）．
双方は canonical compact なので $\RHom(Q_1,X)$ は
perfect で導来基底変更と可換し，fiber は $k[0]$ または零である．
従って matching locus は開閉で，該当側では Hom は可逆層となる．
評価
\[
 Q_1\otimes\Hom(Q_1,X)\longrightarrow X
\]
の perfect cone は全幾何 fiber で零なので，評価は同型である．
非matching側では $Y$ に同じ議論を適用する．
第四項も normal 第四項との perfect Hom と評価，または相対的な
左直交の一意性で同定される．
これで，基底線束 twist を除いて標準四項が実際の normal foundation
となることを確認できる．

従ってこの相対 $H_2$ 族の選択normal窓の有限表示は
実際の $\Gamma_F$，すなわち $\lambda_F$ に一致する．
$\Xi$ の全相対 groupoid 上での全忠実性により，
この canonical cubic tails は $\Xi^{-1}\lambda_F$ と
pointed exact 同値である．
これは fiberwise exact 同値を直接相対化したものではなく，
既に構成済みの二つの相対 $H_2$ 族を，実際の finite normal End を通して
比較している．$U_1$ は $F,G$ の交換で同じである．
\end{proof}

\begin{remark}[用途と範囲]
この命題は二つの開上で，有限 avatar 窓から実際の全次数 cubic 族を
回復する手続きを与える．別に与えた任意の相対 AS/QP 族とこの族が
同型であることは，その族の全整数積と Serre transport を比較して述べる．
また，この段階でも任意の表示 groupoid 全体の同値を主張しない．
\end{remark}
